% =========================================================================
% Figura: Interpretación geométrica del indicador épsilon aditivo (eps+)
% Problema de minimización bidimensional
% =========================================================================

\begin{figure}[htbp]
\centering
\resizebox{0.58\linewidth}{!}{%
\begin{tikzpicture}[
    font=\sffamily,
    >=Stealth,
    x=1cm, y=1cm
]

% -------------------------------------------------------------------------
% 1. EJES COORDENADOS (Minimización)
% -------------------------------------------------------------------------
\draw[line width=1.1pt, draw=darkText, ->, -{Stealth[length=6pt, width=4pt]}]
    (0, 0) -- (6.8, 0) node[below=5pt, midway, font=\fontsize{9.5}{12}\selectfont, text=darkText] {\textbf{Objetivo $f_1$} (minimizar)};

\draw[line width=1.1pt, draw=darkText, ->, -{Stealth[length=6pt, width=4pt]}]
    (0, 0) -- (0, 5.7);
\node[font=\fontsize{9.5}{12}\selectfont, text=darkText, rotate=90]
    at (-0.65, 2.85) {\textbf{Objetivo $f_2$} (minimizar)};

% -------------------------------------------------------------------------
% 2. FRENTE DE REFERENCIA (Curva suave y puntos de referencia Z en figOrange)
% -------------------------------------------------------------------------
\draw[draw=figOrange!75, dashed, line width=1.0pt]
    (0.9, 4.5) to[out=-70, in=150] (2.0, 2.5) to[out=-30, in=165] (3.8, 1.2) to[out=-15, in=175] (6.0, 0.65)
    node[right=2pt, font=\fontsize{8.5}{10.5}\selectfont\bfseries, text=figOrange!90!black] {$PF_{\text{ref}}$};

% Coordenadas de los puntos de referencia Z
\coordinate (z1) at (1.20, 3.70);
\coordinate (z2) at (2.00, 2.50);
\coordinate (z3) at (3.40, 1.40);
\coordinate (z4) at (5.00, 0.85);

% -------------------------------------------------------------------------
% 3. CONJUNTO DE SOLUCIONES OBTENIDAS A (Círculos azules figBlue)
% Desplazadas un valor eps = 0.65 en ambos objetivos
% -------------------------------------------------------------------------
\coordinate (a1) at (1.85, 4.35);
\coordinate (a2) at (2.65, 3.15);
\coordinate (a3) at (4.05, 2.05);
\coordinate (a4) at (5.65, 1.50);

% Trayectoria discontinua del frente obtenido A
\draw[draw=figBlue!40, dashed, line width=0.8pt]
    (1.5, 4.8) to[out=-70, in=150] (2.65, 3.15) to[out=-30, in=165] (4.45, 1.85) to[out=-15, in=175] (6.2, 1.3);

% Componentes de traslacion épsilon en la solucion central a2 (por debajo de la diagonal)
\draw[draw=grayText!60, dashed, line width=0.6pt]
    (a2) -- (2.65, 2.50) node[midway, right=1pt, font=\fontsize{7}{8.5}\selectfont, text=darkText] {$\epsilon$};
\draw[draw=grayText!60, dashed, line width=0.6pt]
    (2.65, 2.50) -- (z2) node[midway, below=1pt, font=\fontsize{7}{8.5}\selectfont, text=darkText] {$\epsilon$};

% -------------------------------------------------------------------------
% 4. VECTORES DE TRASLACION EPSILON (Flechas bermellon hacia el origen)
% -------------------------------------------------------------------------
\draw[line width=1.1pt, draw=figVermilion, -{Stealth[length=4.5pt, width=3.2pt]}] (a1) -- (z1);
\draw[line width=1.1pt, draw=figVermilion, -{Stealth[length=4.5pt, width=3.2pt]}] (a2) -- (z2);
\draw[line width=1.1pt, draw=figVermilion, -{Stealth[length=4.5pt, width=3.2pt]}] (a3) -- (z3);
\draw[line width=1.1pt, draw=figVermilion, -{Stealth[length=4.5pt, width=3.2pt]}] (a4) -- (z4);

% Dibujar puntos de referencia Z (circulos naranja figOrange)
\newcommand{\drawZEps}[1]{%
    \filldraw[draw=figOrange!90!black, fill=figOrange, line width=0.8pt] (#1) circle (0.10cm);
}
\drawZEps{z1}
\drawZEps{z2}
\drawZEps{z3}
\drawZEps{z4}

% Dibujar soluciones A (circulos azules figBlue)
\newcommand{\drawAEps}[1]{%
    \filldraw[draw=figBlue!90!black, fill=figBlue, line width=0.9pt] (#1) circle (0.12cm);
}
\drawAEps{a1}
\drawAEps{a2}
\drawAEps{a3}
\drawAEps{a4}

% -------------------------------------------------------------------------
% 5. LEYENDA (Esquina superior derecha libre)
% -------------------------------------------------------------------------
\begin{scope}[shift={(2.90, 3.75)}]
    \draw[draw=grayText!40, fill=white, rounded corners=3pt, line width=0.6pt]
        (0, 0) rectangle (3.75, 1.30);

    % Soluciones obtenidas A
    \filldraw[draw=figBlue!90!black, fill=figBlue, line width=0.6pt] (0.20, 1.00) circle (0.08cm);
    \node[anchor=west, font=\fontsize{6.5}{7.8}\selectfont, text=darkText]
        at (0.36, 1.00) {Soluciones obtenidas ($A$)};

    % Puntos de referencia Z
    \filldraw[draw=figOrange!90!black, fill=figOrange, line width=0.6pt] (0.20, 0.65) circle (0.07cm);
    \node[anchor=west, font=\fontsize{6.5}{7.8}\selectfont, text=darkText]
        at (0.36, 0.65) {Puntos de referencia ($Z$)};

    % Traslacion épsilon
    \draw[line width=1.0pt, draw=figVermilion, -{Stealth[length=3.5pt, width=2.4pt]}]
        (0.08, 0.30) -- (0.32, 0.30);
    \node[anchor=west, font=\fontsize{6.5}{7.8}\selectfont, text=darkText]
        at (0.36, 0.30) {Traslación ($\epsilon$ aditivo)};
\end{scope}

\end{tikzpicture}%
}
\caption{Interpretación geométrica del indicador $\epsilon$ aditivo en un espacio de dos objetivos a minimizar. Representa el factor mínimo de traslación $\epsilon$ en cada objetivo necesario para que el conjunto obtenido $A$ llegue a dominar débilmente al conjunto de referencia $Z$.}
\label{fig:indicador_epsilon}
\end{figure}
